\chapter{Introduction}

\section{Introduction}
Rice, edible oil, sugar, onion, eggs and lentils take a large share of what a Bangladeshi household spends on food, so a sudden rise in any of them is noticed quickly and often becomes a policy question. The Trading Corporation of Bangladesh (TCB) publishes a retail price report for Dhaka on most working days. The reports are Excel sheets with a minimum and maximum price per item, and they are not available as one clean, machine-readable series.

TCB prices are harder to forecast than they look. On most days a price does not move at all, and then it jumps. For administered items such as sugar and edible oil the jump follows an announcement, while perishables such as onion and eggs move in smaller and more frequent steps. The median item had an unchanged price on 85 to 93 percent of report days, depending on its group (Table~\ref{tab:strata}). So a forecast of ``no change'' is right most of the time, and a model is only worth running if it does better than that.

\section{Motivation}
Two kinds of users could use a daily price watch. A person who follows prices wants to know what a staple is likely to cost in a week or a month and how far from that figure a price could reasonably land. A person who follows market behaviour wants to know when a price has risen further than its own recent history makes likely, and whether the cost drivers that can be measured account for that rise. There is little published forecasting work on Bangladeshi retail prices, and most of it uses price history alone \cite{agripricebd}. I did not find any system that gives the forecast, the range, the alert and the breakdown together for the TCB daily reports.

\section{Problem Statement}
Given the TCB daily retail prices of an item up to a report day, the task is to forecast its price 1, 7, 14 and 30 report days ahead, to give an expected range around each forecast, to decide whether a rise has stayed outside that range long enough to matter, and to show how much of the forecast comes from each driver that can be measured. Because most prices do not move on most days, a model has to beat the forecast ``tomorrow equals today''. A small error on its own proves little.

\section{Objectives of the Study}
The project had five objectives:
\begin{itemize}
    \item to build a clean, audited daily price panel from the TCB reports and join to it the drivers that could be measured, each dated by the day it became known;
    \item to forecast every long-history item at four horizons with one model and to compare that model with five alternatives on a validation year;
    \item to give each forecast an expected range and to raise an alert when the price stays above it, with the alert level fixed from false alarms on non-event days;
    \item to attach a breakdown by driver group to each alert, and to word any residual rise as \unexpl;
    \item to test the final settings once, on a period not used for any choice, and to report the result even if it is unfavourable.
\end{itemize}

\section{Scope of the Study}
The study covers the TCB Dhaka retail reports from 17 May 2021 to 24 October 2026. Of the 63 items, 54 have enough history to forecast. The other nine were added to the report from May 2024 onwards or were reported only for a short time, so they are not forecast. Time is counted in report days, because TCB does not publish every day. The 2022 reports are too sparse and are never filled or scored. Forecasting is for the Dhaka retail price only. The study does not try to find the cause of any price rise. It also does not use market prices from other sources or test pretrained forecasting models. The alerts pick up an abnormal rise after it has started, so they should not be read as an early warning.

\section{Functionality of the Proposed System}
The system reads the price file and produces three things for each of the 54 items. The first is a forecast at 1, 7, 14 and 30 report days with a lower and an upper bound. The second is a list of alerts; an alert starts when the price is above the upper bound on at least 3 of the last 5 report days. The third is, for each alert day, the contribution of each driver group to the forecast and the part of the rise above the range, which is labelled \unexpl. The results go into three files (forecasts, alerts and alert drivers), which are read by a Gradio app in a Colab notebook and by a local Streamlit dashboard. When TCB publishes more reports, a longer price file can be added and the same steps run again.

\section{Organization of the Report}
The report has five chapters.

\textbf{Chapter 1: Introduction} describes the background, problem, objectives, scope and the functions of the system.

\textbf{Chapter 2: Literature Review} reviews related work, states the research gap and lists the contributions.

\textbf{Chapter 3: Methodology} describes the data, the model, its training and the evaluation design.

\textbf{Chapter 4: Experiments and Results} gives the setup and the results for selection, tuning, the sealed test, ranges, alerts and drivers.

\textbf{Chapter 5: Conclusion and Future Work} gives the conclusion, the limitations and the future work.
